% !TeX root = ../main.tex
% Add the above to each chapter to make compiling the PDF easier in some editors.

\chapter{Conclusion}\label{chapter:conclusion}

\section{Summary of Findings}\label{sec:summary-findings}

This thesis investigated the use of convolutional networks for iterative solver selection on sparse linear systems.
The most important finding for the image representation is that the magnitude encoding consistently provides the best information
for the CNN branch, outperforming all other encoding strategies tested across all experiments.
The 14-feature set proved sufficient for the task,
as the six additional features yielded only marginal accuracy gains. This suggests that the CNN image already captures much of the information
the extra features encode.
The MM-AutoSolver architecture was successfully replicated on an independently constructed dataset, confirming the viability of the multimodal approach and
validating their results~(\autoref{sec:mm-replication}).
Extending the single CNN channel to a dual-channel approach, where complementary image encodings are processed in parallel, consistently improved both accuracy and F1 over the
single-channel approach.
An ensemble of four dual-channel models improved the results further, showing that independently trained models with different image mode combinations capture complementary information.
Furthermore, the convergence penalty was found to cause class avoidance rather than real failure rate reduction in some cases, with potential
improvements discussed in \autoref{sec:results-features}.
Overall, the thesis demonstrates that a multimodal approach can meaningfully improve solver selection over non-learning baselines,
given sufficient training data and compute.

\section{Contributions}\label{sec:final-contributions}

\begin{itemize}
  \item Replication of the MM-AutoSolver architecture on an independently
        constructed dataset, validating the multimodal approach.
  \item Design and evaluation of a custom residual CNN architecture in two
        model sizes, supporting both single- and dual-channel image input.
  \item Systematic comparison of 14 image downsampling strategies across both
        64\,px and 128\,px resolutions, identifying magnitude encoding as the
        most effective representation.
  \item Dual-channel CNN evaluation across 30 image mode combinations via a
        shared multimode HDF5 rendering pipeline, demonstrating consistent
        improvement over single-channel models.
  \item Ensemble evaluation of four dual-channel models, showing further
        accuracy and F1 gains through complementary model combination.
  \item Introduction and evaluation of 6 additional matrix features, showing
        that the original 14-feature set is largely sufficient.
  \item Analysis of the convergence penalty loss term, identifying a
        class-avoidance side-effect and suggesting directions for improvement.
  \item Evaluation of inference cost, confirming that prediction overhead
        remains negligible compared to typical iterative solver runtimes.
\end{itemize}

\section{Limitations}\label{sec:limitations}

The dataset, while sufficient for the experiments conducted, is limited in size and diversity.
The synthetic data distribution might not cover all real-world matrix types, and a larger variety of other sources for real-world matrices would be needed
to validate the findings more broadly.
Furthermore, only 19 fixed solver-preconditioner pairs from PETSc were evaluated, even though the full library has many more combinations to offer.
Hyperparameter search was limited to fixed learning rates and architecture depths, leaving potential performance improvements unexplored.
The accuracy drop observed in the \texttt{density\_64} model under the convergence penalty was not fully investigated, and only hypotheses on the cause were discussed.
Finally, the current approach uses hard labels based on the fastest converging solver; soft labels derived from runtime ratios could provide richer training signals.

\section{Future Work}\label{sec:future-work}

\begin{itemize}
  \item Investigation of the \texttt{density\_64} side-effect under the convergence penalty,
        for example through a lambda sweep or mode-specific penalty weights.
  \item An improved convergence penalty formulation, such as the squared penalty or a
        runtime-ratio-based penalty as proposed in \autoref{sec:results-features}.
  \item Soft-label training, where targets are derived from runtime ratios rather than
        a hard best-solver label, to provide richer training signal.
  \item A larger and more diverse dataset, incorporating additional SuiteSparse matrices
        and a wider variety of synthetic matrix types, to improve generalisation.
  \item Extension to larger image resolutions (e.g.\ 256\,px) or learned image encodings
        to capture finer structural detail.
  \item A generalisation study training on synthetic data and evaluating on pure
        SuiteSparse matrices to quantify the distribution shift.
  \item Online integration of the model as a lightweight inference module within the
        embedding-based evaluation pipeline~\cite{weng2026embedding}, to measure real-world solver selection benefit.
  \item Active learning strategies that query the solver oracle only for the most
        uncertain matrices, reducing data generation cost.
  \item Feature selection or pruning to identify a minimal subset of matrix statistics
        that retains high prediction accuracy at reduced extraction cost, informed by the
        per-feature timing analysis in \autoref{sec:inference-cost}.
\end{itemize}
